\documentclass[12pt]{article}

% هوامش
\usepackage{geometry}
\geometry{margin=0.8in}

% ألوان/رؤوس/تذييل/تباعد
\usepackage{xcolor}
\usepackage{fancyhdr}
\usepackage{setspace}
\usepackage{enumitem}

% دعم يونيكود ولغات (يستلزم XeLaTeX)
\usepackage{fontspec}
\usepackage{polyglossia}
\usepackage{bidi} % <<< أضِف هذا السطر

% Math packages for bmatrix/pmatrix/align etc.
\usepackage{amsmath, amssymb, mathtools}

% لصفّ الخيارات أفقياً مع لفّ تلقائي
\usepackage{tabularx, ragged2e, array}
\newcolumntype{Y}{>{\RaggedRight\arraybackslash}X}

% ماكرو جدول الاختيارات (A,B / C,D / E)
\newcommand{\renderchoices}[5]{%
  \begin{tabularx}{\linewidth}{@{}>{\bfseries}p{1.2em}Y >{\bfseries}p{1.2em}Y@{}}
    A.& #1 & B.& #2\\[0.25em]
    C.& #3 & D.& #4\\[0.25em]
    \multicolumn{1}{@{}>{\bfseries}p{1.2em}}{E.}& #5 & \multicolumn{2}{l}{}
  \end{tabularx}\par
}

\setmainlanguage[numerals=western]{arabic} % أرقام 0-9
\setotherlanguage{english}

% الخطوط
\newfontfamily\arabicfont[Script=Arabic,Scale=1.05]{Amiri}
\newfontfamily\englishfont{Times New Roman}

% ترويسة/تذييل
\pagestyle{fancy}
\fancyhf{}
\renewcommand{\headrulewidth}{0pt}
\cfoot{\textenglish{Page } \thepage}

% ماكرو الخيار (اختياري)
\newcommand{\choice}[2]{\textbf{#1}.~#2}

% فقرات
\setlength{\parskip}{0.4em}
\setlength{\parindent}{0pt}

\begin{document}
% ======= Header (fixed layout) =======
\noindent
\begin{tabular*}{\textwidth}{@{\extracolsep{\fill}} l c r}
\textenglish{Date: \textLR{{\textLR{\textenglish{September 4, 2025}}}}} &
{\Large \textbf{{\textLR{\textenglish{Algorithms -1-}}}}} &
\textLR{{\textLR{\textenglish{University of Aleppo}}}} \\
\textenglish{Student Name:} &
\textenglish{Model: \textLR{{\textLR{\textenglish{A}}}}} &
\textenglish{\textbf{Duration:} \textLR{{\textLR{\textenglish{90 minuts}}}}} \\
\textenglish{Student Number:} & & \\
\end{tabular*}

\vspace{0.9cm}

\section*{اختر الإجابة الصحيحة لكل سؤال}
\setstretch{1.25}

\noindent \textbf{\textLR{\textenglish{1.}}} {\textarabic{احسب مشتق التابع }}\( f(x) = x^2\){\textarabic{ وأوجد نهايته عند اللانهاية }}\(\lim_{x \to \infty}f(x)\)

\renderchoices{\textLR{{\textLR{\textenglish{14}}}}}{\textLR{{\textLR{\textenglish{15}}}}}{\underline{\textLR{{\textLR{\textenglish{16}}}}}}{\textLR{{\textLR{\textenglish{17}}}}}{\textLR{{\textLR{\textenglish{17}}}}}
\hfill {\textLR{\textenglish{[1 Point]}}}\\[0.5em]

\noindent \textbf{\textLR{\textenglish{2.}}} \textLR{{\textLR{\textenglish{Please calculate }}}\(\int f(x)dx\)}

\renderchoices{\textLR{{\textLR{\textenglish{13}}}}}{\underline{\textLR{{\textLR{\textenglish{14}}}}}}{\textLR{{\textLR{\textenglish{15}}}}}{\textLR{{\textLR{\textenglish{16}}}}}{}
\hfill {\textLR{\textenglish{[2 Points]}}}\\[0.5em]

\noindent \textbf{\textLR{\textenglish{3.}}} \textLR{{\textLR{\textenglish{Test}}}}

\renderchoices{\textLR{{\textLR{\textenglish{13}}}}}{\underline{\textLR{{\textLR{\textenglish{14}}}}}}{\textLR{{\textLR{\textenglish{15}}}}}{\textLR{{\textLR{\textenglish{16}}}}}{\textLR{{\textLR{\textenglish{17}}}}}
\hfill {\textLR{\textenglish{[3 Points]}}}\\[0.5em]

\noindent \textbf{\textLR{\textenglish{4.}}} \textLR{{\textLR{\textenglish{Test2}}}}

\renderchoices{\underline{\textLR{{\textLR{\textenglish{1}}}}}}{\textLR{{\textLR{\textenglish{2}}}}}{\textLR{{\textLR{\textenglish{3}}}}}{\textLR{{\textLR{\textenglish{4}}}}}{\textLR{{\textLR{\textenglish{5}}}}}
\hfill {\textLR{\textenglish{[3 Points]}}}\\[0.5em]

\noindent \textbf{\textLR{\textenglish{5.}}} \textLR{{\textLR{\textenglish{Test3}}}}

\renderchoices{\textLR{{\textLR{\textenglish{1}}}}}{\textLR{{\textLR{\textenglish{2}}}}}{\underline{\textLR{{\textLR{\textenglish{3}}}}}}{\textLR{{\textLR{\textenglish{4}}}}}{\textLR{{\textLR{\textenglish{5}}}}}
\hfill {\textLR{\textenglish{[1 Point]}}}\\[0.5em]

\noindent \textbf{\textLR{\textenglish{6.}}} \textLR{{\textLR{\textenglish{test4}}}}

\renderchoices{\underline{\textLR{{\textLR{\textenglish{1}}}}}}{\textLR{{\textLR{\textenglish{2}}}}}{\textLR{{\textLR{\textenglish{3}}}}}{\textLR{{\textLR{\textenglish{4}}}}}{\textLR{{\textLR{\textenglish{5}}}}}
\hfill {\textLR{\textenglish{[1 Point]}}}\\[0.5em]

\noindent \textbf{\textLR{\textenglish{7.}}} \textLR{{\textLR{\textenglish{Hello!}}}}

\renderchoices{\underline{\textLR{{\textLR{\textenglish{1}}}}}}{\textLR{{\textLR{\textenglish{2}}}}}{\textLR{{\textLR{\textenglish{3}}}}}{\textLR{{\textLR{\textenglish{4}}}}}{}
\hfill {\textLR{\textenglish{[2 Points]}}}\\[0.5em]

\noindent \textbf{\textLR{\textenglish{8.}}} \textLR{{\textLR{\textenglish{Find the inverse of the following matrix: }}}\(\begin{bmatrix} 1 & 3 & 4 \\ 5 & 2 & 1 \\ 6 & 8 & 1 \end{bmatrix}\)}

\renderchoices{\underline{\textLR{{\textLR{\textenglish{There is no inverse for this matrix}}}}}}{\textLR{\(\begin{pmatrix} 1 & 2 & 3 \\ 4 & 5 & 6 \\ 7 & 8 & 9  \end{pmatrix}\)}}{\textLR{{\textLR{\textenglish{0}}}}}{\textLR{{\textLR{\textenglish{0}}}}}{}
\hfill {\textLR{\textenglish{[4 Points]}}}\\[0.5em]

\noindent \textbf{\textLR{\textenglish{9.}}} \textLR{{\textLR{\textenglish{Solve for }}}\(x \){\textLR{\textenglish{ in the equation }}}\(2x+6 = 0 \)}

\renderchoices{\textLR{{\textLR{\textenglish{-1}}}}}{\textLR{{\textLR{\textenglish{-2}}}}}{\underline{\textLR{{\textLR{\textenglish{-3}}}}}}{\textLR{{\textLR{\textenglish{3}}}}}{}
\hfill {\textLR{\textenglish{[1 Point]}}}\\[0.5em]

\noindent \textbf{\textLR{\textenglish{10.}}} \textLR{\(x=3+6 \){\textLR{\textenglish{ what is the value of x?}}}}

\renderchoices{\textLR{{\textLR{\textenglish{7}}}}}{\underline{\textLR{{\textLR{\textenglish{9}}}}}}{\textLR{{\textLR{\textenglish{8}}}}}{\textLR{{\textLR{\textenglish{10}}}}}{\textLR{{\textLR{\textenglish{0}}}}}
\hfill {\textLR{\textenglish{[1 Point]}}}\\[0.5em]



\end{document}
